% Part: model-theory
% Chapter: interpolation
% Section: interpolation-proof

\documentclass[../../../include/open-logic-section]{subfiles}

\begin{document}

\olfileid{mod}{int}{prf}

\olsection{کریگ دا درمیانی جملے دا قضیہ}

\begin{thm}[کریگ دا درمیانی جملے دا قضیہ]
\ollabel{thm:interpol} 
فرض کرو دتے دو فارمولے جملے نیں تے
$\Entails !A \lif !B$، تاں اک !!a{sentence}
$!C$ موجود اے جس لئی
$\Entails !A \lif !C$ تے
$\Entails !C \lif !B$، تے ہر
!!{constant}، !!{function} تے
!!{predicate} ($\eq$ توں علاوہ)
جو $!C$ وچ اے، $!A$
تے~$!B$ دوہاں وچ وی آؤندا اے۔
!!{sentence} $!C$ نوں $!A$ تے~$!B$
دا \emph{درمیانی جملہ (interpolant)}
آکھیا جاندا اے۔
\end{thm}

\begin{proof}
فرض کرو $\Lang{L}_1$، $!A$ دی زبان،
تے $\Lang{L}_2$، $!B$ دی زبان اے۔
$\Lang{L}_0 = \Lang{L}_1 \cap \Lang{L}_2$
رکھو۔ ہر $i \in \{0, 1, 2 \}$ لئی
$\Lang{L}'_i$، $\Lang{L}_i$ وچ
لامتناہی نویں !!{constant}s $\Obj c_0,
\Obj c_1, \Obj c_2, \dots$ شامل کرکے بناؤ۔

جے $!A$ پورا ہون جوگا نہ ہووے،
$\lexists[x][\eq/[x][x]]$ اک درمیانی
جملہ اے۔ جے $\lnot !B$ پورا ہون جوگا
نہ ہووے (تے ایس لئی $!B$ منطقی طور
اُتے معتبر ہووے)، تاں
$\lexists[x][\eq[x][x]]$ اک درمیانی
جملہ اے۔ سو اسیں ایہ وی فرض کر سکدے آں
کہ $!A$ تے $\lnot !B$ دوہاں
پورے ہون جوگے نیں۔

درمیانی جملے دے قضیے دا تقابلِ نقیض
ثابت کرن لئی فرض کرو کہ $!A$ تے
$!B$ لئی کوئی درمیانی جملہ نہیں۔
دوجے لفظاں وچ فرض کرو کہ $\{!A\}$
تے $\{\lnot !B\}$،
$\Lang{L}_0$ وچ وکھ نہ کیتے
جا سکدے نیں۔

ساڈا مقصد جوڑی $(\{ !A \}, \{\lnot!B\})$
نوں اک ایسی غیر قابلِ تفریق جوڑی
$(\Gamma^*, \Delta^*)$ تک ودھانا اے
جس نوں ہور نہ ودھایا جا سکے۔
$!A_0$، $!A_1$، $!A_2$، \dots،
$\Lang{L}'_1$ دے !!{sentence}s دی
گنتی وار فہرست ہون؛ تے $!B_0$،
$!B_1$، $!B_2$، \dots،
~$\Lang{L}'_2$ دے !!{sentence}s دی۔
اسیں !!{sentence}s دے سیٹاں دیاں دو
ودھدیاں سلسلیاں $(\Gamma_n, \Delta_n)$،
$n \ge 0$ لئی، ایویں متعین کردے آں۔
$\Gamma_0 = \{ !A\}$ تے
$\Delta_0 = \{\lnot !B \}$ رکھو۔
فرض کرو $(\Gamma_n, \Delta_n)$ پہلے
توں متعین نیں؛ $\Gamma_{n+1}$ تے
$\Delta_{n+1}$ ایویں متعین کرو:
ہر اگلا سیٹ پچھلا سیٹ رکھدا اے؛ ہیٹھاں دتی شرط
پوری ہووے تاں صرف دتے اضافے کرو، نہیں تاں اوہ سیٹ
جیویں اے رہن دیو۔ گواہ لئی شامل کیتے مشترک نویں
مستقلاں وچوں اک چنو:
\begin{enumerate}
\item جے $\Gamma_n \cup \{!A_n \}$ تے
  $\Delta_n$، $\Lang{L}'_0$ وچ
  وکھ نہ کیتے جا سکّن، تاں $!A_n$ نوں
  $\Gamma_{n+1}$ وچ رکھو۔ ہور ایہ کہ جے
  $!A_n$ وجودی !!{formula}
  $\lexists[x][!S]$ ہووے، تاں ایسا نواں
  !!{constant} $c$ چنو جو
  $\Gamma_n$، $\Delta_n$، $!A_n$ یا
  $!B_n$ وچ نہ آؤندا ہووے، تے
  $\Subst{!S}{c}{x}$ نوں
  $\Gamma_{n+1}$ وچ رکھو۔
\item جے $\Gamma_{n+1}$ تے
  $\Delta_n \cup \{!B_n \}$،
  $\Lang{L}'_0$ وچ وکھ نہ کیتے جا سکّن،
  تاں $!B_n$ نوں $\Delta_{n+1}$ وچ رکھو۔
  ہور ایہ کہ جے $!B_n$ وجودی
  !!{formula} $\lexists[x][!S]$ ہووے،
  تاں ایسا نواں !!{constant} $c$
  چنو جو $\Gamma_{n+1}$، $\Delta_n$،
  $!A_n$ یا $!B_n$ وچ نہ آؤندا ہووے،
  تے $\Subst{!S}{c}{x}$ نوں
  $\Delta_{n+1}$ وچ رکھو۔
\end{enumerate}
آخر وچ ایہ متعین کرو:
\begin{align*}
  \Gamma^* & = \bigcup_{n\ge 0} \Gamma_n, & 
  \Delta^* & = \bigcup_{n\ge 0} \Delta_n.
\end{align*}
ہن $n$ اُتے دوہری استقرا نال
اسیں ثابت کر سکدے آں:
\begin{enumerate}
\item\ollabel{part-a} $\Gamma_n$ تے
  $\Delta_n$، $\Lang{L}'_0$
  وچ وکھ نہ کیتے جا سکّن؛
\item\ollabel{part-b} $\Gamma_{n+1}$
  تے $\Delta_n$، $\Lang{L}'_0$
  وچ وکھ نہ کیتے جا سکّن۔
\end{enumerate}
\olref{part-a} دی ابتدا
\olref[sep]{lem:sep1} توں ملدی اے۔
\olref{part-b} لئی تین صورتاں
وکھ کرنیاں پین گیاں:
\begin{enumerate}
\item جے $\Gamma_0 \cup \{!A_0 \}$
  تے $\Delta_0$ وکھ کیتے جا سکّن،
  تاں $\Gamma_1 = \Gamma_0$ تے
  \olref{part-b} صرف \olref{part-a}
  اے؛
\item جے $\Gamma_1 = \Gamma_0 \cup\{ !A_0\}$،
  تاں بناوٹ دے مطابق $\Gamma_1$ تے
  $\Delta_0$ وکھ نہیں کیتے جا سکدے۔
\item ہن اوہ صورت رہندی اے جدوں $!A_0$
  وجودی ہووے، سو
  $\Gamma_1 = \Gamma_0 \cup \{ \lexists[x][!S], \Subst{!S}{c}{x}
  \}$۔ بناوٹ دے مطابق
  $\Gamma_0 \cup \{ \lexists[x][!S]\}$
  تے $\Delta_0$ وکھ نہیں کیتے جا سکدے؛
  سو \olref[sep]{lem:sep2} نال
  $\Gamma_0 \cup \{ \lexists[x][!S], \Subst{!S}{c}{x} \}$
  تے $\Delta_0$ وی وکھ نہیں کیتے جا سکدے۔
\end{enumerate}
ایس نال اُتے \olref{part-a} تے
\olref{part-b} لئی استقرا دی ابتدا
مکمل ہوندی اے۔ ہن استقرائی قدم۔
\olref{part-a} لئی، جے
$\Delta_{n+1} \supseteq \Delta_n \cup \{ !B_n \}$
تاں $\Gamma_{n+1}$ تے
$\Delta_{n+1}$ بناوٹ دے مطابق وکھ
نہیں کیتے جا سکدے (جے $!B_n$
وجودی وی ہووے تاں
\olref[sep]{lem:sep2} نال)؛ جے
$\Delta_{n+1} = \Delta_n$ ہووے
(کیوں جے $\Gamma_{n+1}$ تے
$\Delta_n \cup \{!B_n\}$ وکھ کیتے
جا سکدے نیں)، تاں
\olref{part-b} دا استقرائی مفروضہ ورتدے
آں۔ \olref{part-b} دے استقرائی قدم لئی،
جے $\Gamma_{n+2} \supseteq \Gamma_{n+1} \cup
\{!A_{n+1} \}$ تاں $\Gamma_{n+2}$
تے $\Delta_{n+1}$ بناوٹ دے مطابق
وکھ نہیں کیتے جا سکدے (جے $!A_{n+1}$
وجودی وی ہووے تاں
\olref[sep]{lem:sep2} نال)؛ تے جے
$\Gamma_{n+2} = \Gamma_{n+1}$ ہووے،
تاں \olref{part-a} دا ہن ثابت ہویا
استقرائی قدم ورتدے آں۔ ایس نال
\olref{part-a} تے \olref{part-b}
اُتے استقرا مکمل ہوندی اے۔

ایس توں $\Gamma^*$ تے $\Delta^*$
وکھ نہیں کیتے جا سکدے؛ جے ایسا نہ
ہووے، تاں محدود گواہی دے قضیے نال
کوئی $n \ge 0$ ایسا ہووے گا کہ
$\Gamma_n$ تے $\Delta_n$ نوں
وکھ کیتا جا سکے، \olref{part-a} دے خلاف۔
خاص طور اُتے $\Gamma^*$ تے
$\Delta^*$ بے تضاد نیں: کیوں جے جے
پہلا یا دوجا سیٹ متضاد ہووے، تاں
اوہناں نوں بالترتیب
$\lexists[x][\eq/[x][x]]$ یا
  $\lforall[x][\eq[x][x]]$
نال وکھ کیتا جا سکدا اے۔

ہن اسیں دکھاؤندے آں کہ $\Gamma^*$،
$\Lang{L}'_1$ وچ مکمل بے تضاد اے
تے ایویں $\Delta^*$،
$\Lang{L}'_2$ وچ۔ پہلے لئی فرض کرو
کہ $!A_n \notin \Gamma^*$ تے
$\lnot !A_n
\notin \Gamma^*$، کسے
$n \ge 0$ لئی۔ جے
$!A_n \notin \Gamma^*$،
تاں $\Gamma_n \cup \{!A_n \}$
نوں $\Delta_n$ توں وکھ کیتا
جا سکدا اے، تے ایس لئی کوئی
$!C \in \Lang{L}'_0$ اے جس لئی
دوہاں گلاں ایہ نیں:
\begin{align*}
  \Gamma^* & \Entails !A_n \lif !C, & 
  \Delta^* & \Entails \lnot !C.
\end{align*}
ایویں جے $\lnot !A_n \notin \Gamma^*$،
تاں کوئی $!C' \in
\Lang{L}'_0$ اے جس لئی دوہاں گلاں:
\begin{align*}
  \Gamma^* & \Entails \lnot !A_n \lif !C', & 
  \Delta^* & \Entails \lnot !C'.
\end{align*}
قضیاتی منطق نال
$\Gamma^* \Entails !C \lor !C'$ تے
$\Delta^* \Entails \lnot (!C \lor !C')$؛
سو $!C \lor
!C'$، $\Gamma^*$ تے $\Delta^*$
نوں وکھ کردا اے۔ ایس ورگی دلیل
دکھاؤندی اے کہ $\Delta^*$ وی مکمل اے۔

آخر وچ اسیں دکھاؤندے آں کہ
$\Gamma^* \cap \Delta^*$،
$\Lang{L}'_0$ وچ مکمل بے تضاد اے۔
ایہ ظاہر طور اُتے بے تضاد اے، کیوں
جے ایہ بے تضاد سیٹاں دا مشترک حصہ اے۔
مکمل پن دکھاؤن لئی $!S \in
\Lang{L}'_0$ نوں اک جملہ من لو۔ ہن $\Gamma^*$،
$\Lang{L'_1}
\supseteq \Lang{L'_0}$ وچ مکمل اے،
تے ایویں $\Delta^*$،
$\Lang{L'_2} \supseteq \Lang{L'_0}$
وچ۔ ایس لئی یا $!S \in \Gamma^*$
یا $\lnot !S \in \Gamma^*$، تے
یا $!S \in \Delta^*$ یا
$\lnot !S \in \Delta^*$۔ جے
$!S \in
\Gamma^*$ تے $\lnot !S \in \Delta^*$
تاں $!S$، $\Gamma^*$ تے
$\Delta^*$ نوں وکھ کردا؛ تے جے
$\lnot !S \in
\Gamma^*$ تے $!S \in \Delta^*$
تاں $\Gamma^*$ تے $\Delta^*$
نوں $\lnot !S$ وکھ کردا۔ ایس لئی
یا $!S \in
\Gamma^* \cap \Delta^*$ یا
$\lnot !S \in \Gamma^* \cap \Delta^*$،
تے $\Gamma^* \cap \Delta^*$
مکمل اے۔

بناوٹ وچ ہر قبول کیتے وجودی جملے نال گواہ شامل کیتا گیا اے۔
ہر بند ٹرم دے برابر ہون والا عنصر موجود ہون دی شرط ہمیشہ
صادق اے؛ ودھائی زبان دی فہرست وچ ایہ شرط وی آؤندی اے
تے اوہدے لئی شامل کیتا مشترک نواں مستقل گواہ ملدا اے۔
سو بند ٹرماں دے برابری والے ہر طبقے دا نمائندہ ایسا مستقل اے۔
کیوں جے $\Gamma^*$ مکمل بے تضاد تے گواہاں والا اے،
ایس دا اک نمونہ $\Struct{M}'_1$ اے
جس دا !!{domain} $\Domain{M'_1}$
عین اوہناں عنصراں اُتے مشتمل اے
جو شامل کیتے مشترک نویں !!{constant}s دی تعبیر
$\Assign{c}{M'_1}$ نال ملدے نیں؛
بالکل تمامیت دے قضیے دے ثبوت وانگ
(\olref[fol][com][cth]{thm:completeness})۔
ایویں $\Delta^*$ دا اک نمونہ
$\Struct{M}'_2$ اے جس دا !!{domain}
$\Domain{M'_2}$،
شامل کیتے مشترک نویں !!{constant}s دیاں تعبیراں
$\Assign{c}{M'_2}$ نال ملدا اے۔

$\Struct{M_1}$، $\Struct{M'_1}$
توں $\Lang{L'_1} \setminus \Lang{L'_0}$
وچلی !!{constant}s، !!{function}s تے
!!{predicate}s دیاں تعبیراں ہٹا کے
بناؤ؛ تے $\Struct{M_2}$ وی ایسے طرح۔
ایتھے نقشے دی تعریف وچ مستقل شامل کیتے مشترک نویں
مستقلاں وچوں ورتیا جاندا اے؛ پچھلے قدم نال ایہ دوہاں
دائریاں دے ہر عنصر نوں ناں دیندے نیں۔
تاں نقشہ $h \colon M_1 \to M_2$
جو $h(\Assign{c}{M'_1}) = \Assign{c}{M'_2}$
نال متعین اے، $\Lang{L}'_0$
وچ ہم شکلی اے، کیوں جے
$\Gamma^* \cap \Delta^*$،
$\Lang{L}'_0$ وچ مکمل بے تضاد اے،
جیویں دکھایا۔ ایہ ایس لئی نکلدا
اے کہ کوئی وی $\Lang{L}'_0$-!!{sentence}
یا $\Gamma^*$ تے $\Delta^*$ دوہاں
وچ اے یا کسے وچ وی نہیں: سو
$\Assign{c}{M'_1} \in
\Assign{P}{M'_1}$ اوہدوں تے صرف
اوہدوں جدوں $\Atom{P}{c} \in \Gamma^*$،
اوہدوں تے صرف اوہدوں جدوں
$\Atom{P}{c} \in \Delta^*$، اوہدوں
تے صرف اوہدوں جدوں
$\Assign{c}{M'_2} \in
\Assign{P}{M'_2}$۔ ہم شکلی دیاں
باقی شرطاں وی ایویں ثابت کیتیاں
جا سکدیاں نیں۔

ہن اک نمونہ $\Struct{M}$،
!!{language}
$\Lang{L'_1} \cup \Lang{L'_2}$ لئی
ایویں متعین کریے:
\begin{enumerate}
\item !!{domain} $\Domain{M}$،
  عین $\Domain{M_2}$ اے، یعنی سارے
  عنصر $\Assign{c}{M'_2}$ دا سیٹ؛
\item جے !!a{predicate}~$P$،
  $\Lang{L_2} \setminus
  \Lang{L_1}$ وچ اے، تاں
  $\Assign{P}{M} = \Assign{P}{M'_2}$؛
\item جے محمول $P$،
  $\Lang{L}_1\setminus \Lang{L}_2$
  وچ اے، تاں
  $\Assign{P}{M} = h(\Assign{P}{M'_1})$،
  یعنی $\tuple{\Assign{c_1}{M'_2}, \dots, \Assign{c_n}{M'_2}} \in
  \Assign{P}{M}$ اوہدوں تے صرف اوہدوں
  جدوں $\tuple{\Assign{c_1}{M'_1}, \dots,
  \Assign{c_n}{M'_1}} \in \Assign{P}{M'_1}$۔
\item جے !!a{predicate} $P$،
  $\Lang{L}_0$ وچ اے، تاں
  $\Assign{P}{M} =
  \Assign{P}{M'_2} = h(\Assign{P}{M'_1})$۔
\item $\Lang{L}'_1 \cup \Lang{L}'_2$
  دے !!^{function}s، جنہاں وچ
  !!{constant}s وی نیں، ایویں ہی
  سمجھے جاندے نیں۔
\end{enumerate}

آخر وچ !!{formula}s اُتے استقرا
نال دکھایا جاندا اے کہ
$\Struct{M}$، $\Struct{M'_1}$
نال $\Lang{L'_1}$ دے سارے
!!{formula}s بارے، تے
$\Struct{M'_2}$ نال
$\Lang{L'_2}$ دے سارے
!!{formula}s بارے متفق اے؛ پہلے مقابلے وچ ہم شکلی
راہیں ملدیاں متغیر قدر-بندیاں، تے دوجے وچ اکو قدر-بندی
ورتی جاندی اے۔ جملیاں لئی قدر-بندی دا فرق نہیں پیندا۔
خاص طور اُتے
$\Struct{M} \Entails \Gamma^* \cup \Delta^*$،
جس توں $\Struct{M}
\Entails !A$ تے
$\Struct{M} \Entails \lnot!B$،
تے $\not\Entails !A \lif !B$۔
ایس نال کریگ دے درمیانی جملے دے
قضیے دا ثبوت مکمل ہوندا اے۔
\end{proof}

\end{document}
